% caption_supp.tex -- GENERATED by figures/v7/supp/scripts/make_captions.py; do not edit by hand.
% Captions of Supplementary Figures S1-S10 (v2). Every number is rendered from analysis/outputs/stats.json or a
% derived_numbers.json by the script (log: panels/captions_numbers.csv). Figure files: figS1.pdf ... figS10.pdf.

% ---- Figure S1 (figS1.pdf)
\caption{\textbf{Shopify's Web Pixels Manager (WPM) decides in the page's main thread whether a sandboxed pixel starts, and never stops it afterwards.} A schematic of our static reading of the WPM manager bundle and its inline loader (Methods), read from left to right and drawn like Fig.~1a: the outer frame is the Shopify store in the visitor's browser; inside it, the light grey ground is the page's main thread and the darker grey ground beside it the pixel sandboxes, which are isolated from the page; the outlined frame inside the page is the manager bundle. Boxes are components, with their secondary text in grey; arrows show what feeds what. Inline loader (top, in the page): each storefront page lists every pixel (`pixel list') with its type (`APP' or `CUSTOM'), its runtime context (`runtimeContext': `STRICT', `LAX' or `OPEN') and its declared privacy purposes (analytics, marketing, preferences, sale of data), so a crawler can read the list before any script runs; the list feeds the consent gate. Consent signals (left), in the order and words of Fig.~1a: the choice of a consent-management platform (CMP) that calls `setTrackingConsent'; the merchant's region default, written by the server into the page response; and, for sale of data only, Global Privacy Control (GPC), whose `Sec-GPC' request header makes the server decline sale of data. Manager bundle (middle): the consent state is the first one found of the sources listed in its box, in this order: `cachedConsent', `injectedConsent', the `\_cs' URL parameter, the `\_tracking\_consent' cookie and the `Server-Timing: \_cmp' header; if none can be read, everything is allowed. The event bus replays page events to every pixel, and each event carries the full page URL. The consent gate starts a pixel only when the consent state allows every purpose the pixel declares, and starts a pixel that declares no purpose at once. Runtime containers (right; the grey qualifier in each box is the runtime context): the gate starts each pixel in the container its runtime context names, a web worker (strict, `STRICT') or a sandboxed iframe (lax, `LAX'), both on the pixel sandboxes' ground, so a script blocker in the page cannot see them, or a script in the page (open, `OPEN'), drawn on the page's ground, which a blocker can see. In the web worker the vendor library sends with `fetch' without cookies, and `navigator' is overwritten, so pixel code reads GPC as undefined. We found no code path that stops a pixel once it has started. Schematic drawn in TikZ, with no generated imagery.}
\label{fig:s1}

% ---- Figure S2 (figS2.pdf)
\caption{\textbf{Shopify's own banner had the lowest pre-choice firing of any consent tool and also the most denying region defaults, so the tools cannot be ranked on these data.} (a)~Six outcomes per consent-management platform (CMP), the tool that showed the banner in the Default load. The rows are grouped by the tool's integration mode, each group set apart by a gap and the four mode groups of third-party tools named by a grey header; from top to bottom: Shopify's own banner; tools that both write the choice to Shopify and block tags in the page's main thread (`Writes and blocks': Pandectes, OneTrust, Cookiebot, Osano, Ketch, Clarip and TrustArc); a tool that writes to Shopify only (`Writes only': iubenda); a main-thread blocker only (`Blocks only': Usercentrics); tools whose mode is unknown (`Mode unknown': Axeptio, CookieScript and CookieYes); and other or no detected tool. One dot plot per outcome, the tools sharing rows; from left to right: Shopify's region default denies marketing (banner stores with a readable default) and a sandboxed ad request before any choice (banner stores), both in the Default load; a reject control on the banner's first layer (stores where the Decline load found a banner), Shopify recorded an explicit refusal (working clicks) and a new sandbox sent an ad request on the fresh page after Decline (working clicks), all three in the Decline load; and a sandboxed ad request still sent under Global Privacy Control (GPC; stores with one under Default, assigned to the tool that showed the banner or else the one detected, so the last row also counts stores without a banner). Marks: share of stores; lines: Wilson 95\% confidence intervals (CIs), in the colour of their mark. Shopify's own banner, the row the title is about, is drawn as slate squares; every other tool is context, drawn as grey circles, and tools with fewer than 10 banner stores are drawn lighter (light grey circles and intervals, names in grey), as the top rank bucket of Fig.~S3. The plots of the three rates that Fig.~S5 tests (fired before any choice, fired on the fresh page and still fired under GPC) carry a pale brick band: that rate over all tools, which is the main estimate of Fig.~S5 with its Wilson 95\% CI (the tools' rows pool to it), so each tool reads against the overall rate. Cells with fewer than 5 stores are not drawn. Shopify's own banner: the region default denied marketing on 65 of 66 stores (98.5\%) and a sandboxed ad request preceded any choice on 12 of 68 (17.6\%). The $k$ and $n$ of every cell are in Table~\ref{tab:s-cmp}. (b)~Share of configured sandboxed vendor pixels that sent an ad request before any choice, by consent tool (rows, in the order of (a)) and vendor group (columns), among banner stores whose Shopify configuration lists the vendor's pixel, in four equal grey steps from light (low) to dark (high; thin colour bar at right), the cells split by thin white gaps; the lightest step is a definite grey, so it never reads as white. The column set apart at the right, `All vendors', counts the store--vendor pairs of every group. By row, the pairs that fired: Shopify's banner 14 of 141 pairs (68 banner stores); Pandectes 124 of 202 (95); OneTrust 76 of 142 (91); Cookiebot 43 of 70 (45); Osano 21 of 46 (17); iubenda 14 of 18 (10); all other tools 99 of 145 (83); in all 391 of 764 pairs (51.2\%). Blank cells with a grey dash: fewer than 10 banner stores configure the vendor's pixel (often none), so no rate is shown; a dash means no estimate, not a low rate. Rates describe the stores that use each tool, not the tool's quality.}
\label{fig:s2}

% ---- Figure S3 (figS3.pdf)
\caption{\textbf{Pre-choice firing on banner stores was about the same in every Tranco rank bucket below the top one.} Outcomes by Tranco rank bucket (rows, top to bottom): the top 10k ranks (7 captured stores), 10k--50k (171), 50k--100k (621) and 100k--200k (559), unweighted. One dot plot per outcome, the buckets sharing rows; from left to right: a banner shown, and a sandboxed ad request before any choice among all stores and among banner stores (Default load); a new sandbox sent an ad request on the fresh page after a working Decline; a sandboxed ad request still sent under Global Privacy Control (GPC; stores with one under Default). The rows are named once, under the grey header `Tranco rank'. Circles: share of stores; lines: Wilson 95\% confidence intervals (CIs), in the colour of their mark. The last three plots show the three rates that Fig.~S5 tests, all rates at which consent was lost, and are drawn alike: brick circles on a pale brick band, which is that rate over all four buckets, the main estimate of Fig.~S5 with its Wilson 95\% CI (the buckets pool to it); the first two plots are context, in grey. The top bucket holds only 7 stores (fewer than 10) and is drawn lighter in every plot (light grey circles and intervals, its name in grey), as the small tools of Fig.~S2a (in Fig.~S5, by contrast, grey marks a changed definition of the measure). Cells with fewer than 5 stores are not drawn. $k$ of $n$ by bucket, top to bottom. Banner shown: 3 of 7; 61 of 171; 161 of 621; 184 of 559. Before any choice, all stores: 6 of 7; 101 of 171; 373 of 621; 323 of 559. Before any choice, banner stores (47.5--50.9\% in the three lower buckets): 2 of 3 (not drawn); 29 of 61; 82 of 161; 92 of 184. Fresh page: 0 of 1 (not drawn); 4 of 32; 20 of 92; 29 of 106. Under GPC: 3 of 6; 34 of 101; 122 of 373; 120 of 323.}
\label{fig:s3}

% ---- Figure S4 (figS4.pdf)
\caption{\textbf{The crawler's banner, reject and click labels agreed with screenshots in almost every audited load, and repeated loads agreed almost perfectly.} (a--c)~Crawler label (rows, shared by the three matrices and named once, left of (a)) against a screenshot audit (columns) in 54 seeded Decline loads (36 clicked, 18 without a reject option), as matrices of loads split into cells by crossing rules, as in Fig.~S7b; the diagonal (top left and bottom right): the crawler and the audit agree; the other two cells: they disagree. (a)~Banner detection: 53 of 54 detected banners were real (the exception was an age gate); the audit sampled only loads with a detected banner, so the crawler's `no' row was not audited (one merged row, marked `not audited'). (b)~Reject option on the banner's first layer: agreement in 52 of 53 loads (one reject button in Romanian was missed). (c)~Whether the Decline click dismissed the banner, in the 35 clicked loads whose screenshot showed a banner: 35 of 35. The audit labels have not yet been checked by an independent human rater. (d)~Cohen's $\kappa$ for whether a sandboxed ad request fired, between a load and its repeat (Default; Global Privacy Control, GPC) and between emulated GPC and Brave's native GPC, per store (filled circles) and per store $\times$ vendor pair (open circles); lines: 95\% confidence intervals from a bootstrap that resamples stores; a comparison without any disagreement has a degenerate interval and no line. Note the truncated $\kappa$ axis, which does not start at zero; it runs a little past one, so the marks at $\kappa = 1$ have air above them. The key sits right of the plot. Default vs repeat: $\kappa$ 0.985 (0.961--1.000), 271 of 273 stores agreeing, and 0.985 (0.967--0.997), 1,902 of 1,911 pairs; GPC vs repeat: 273 of 273 stores ($\kappa$ 1.000) and $\kappa$ 0.982 (0.939--1.000), 1,908 of 1,911 pairs; GPC vs Brave: 98 of 98 stores and 686 of 686 pairs ($\kappa$ 1.000 for both). Default and GPC differed on 39 of the 98 Brave stores and on 109 of 686 store $\times$ vendor pairs, and Brave matched GPC on all of them. Counts of sandboxed ad requests per store also repeated: Spearman's $\rho$ 0.982 (0.957--0.999) for Default and 0.999 (0.995--1.000) for GPC (273 stores; bootstrap over stores).}
\label{fig:s4}

% ---- Figure S5 (figS5.pdf)
\caption{\textbf{Robustness variants moved the three headline rates little; only changes of definition moved them much.} Each row recomputes one headline rate: (a)~banner stores with a sandboxed ad request before any choice (Default load); (b)~working Decline clicks after which a new sandbox sent an ad request on the fresh page; (c)~stores with a sandboxed ad request under Default that still had one under Global Privacy Control (GPC). Filled brick circle: the main estimate with its Wilson 95\% confidence interval (CI), which the pale brick band carries through the plot (a row whose interval reaches into the band overlaps the main interval; the same band marks these rates in Figs~S2a and S3); `Main, bootstrap by consent tool' is the same main estimate with the bias-corrected and accelerated (BCa) cluster-bootstrap interval in place of the Wilson interval; it is marked by a short vertical brick tick on its interval, not a circle (`Main, bootstrap interval' in the key), because it is not a new estimate, only a wider interval around the same one. Open brick circles: robustness variants (crawl segments at 6 and 10 concurrent browsers; edge-cached pages removed; California consent strings only; the 11 banners the detector missed added; the 54 flagged Consentmo stores counted as banner stores; failed forced clicks included, or all forced clicks removed; the first 10~s of each load; stores flagged in the ethics review removed; design weights with a ratio estimator). Grey circles, set apart below a gap: changed definitions of the measure (data beacons only; beacons carrying our test code; any origin; any time after the click; any sandbox on the next page; pixels declaring sale of data only; all GPC stores as the denominator). Lines: 95\% CIs. The rows, top to bottom, as $k$ of $n$ stores: (a)~main 50.1\% (205 of 409; CI 45.3--54.9\%); BCa interval 32.6--58.9\%; variants 19 of 34; 184 of 373; 203 of 404; 204 of 400; 208 of 420; 215 of 463; 201 of 401; weighted 50.1\%; changed definitions 200 of 409; 192 of 409; 357 of 409. (b)~main 22.9\% (53 of 231; CI 18.0--28.8\%); BCa interval 17.4--29.5\%; variants 57 of 240; 7 of 19; 46 of 212; 53 of 228; 52 of 227; weighted 24.9\%; changed definitions 52 of 231; 80 of 231; 70 of 231. (c)~main 34.7\% (279 of 803; CI 31.5--38.1\%); BCa interval 30.9--38.7\%; variants 279 of 802; 20 of 70; 256 of 730; 276 of 798; 275 of 793; weighted 35.8\%; changed definitions 242 of 769; 229 of 753; 0 of 803; 1,213 of 1,254; 279 of 1,354 (the version of this rate reported in the main text). All 20 robustness variants overlap the main interval, but most are near-complete subsets of the main sample, and the 6-browser segment could detect only differences of about 17--30 percentage points (80\% power).}
\label{fig:s5}

% ---- Figure S6 (figS6.pdf)
\caption{\textbf{Before any choice, ad pixels of every vendor group fired from Shopify's sandboxes, and most first sandboxed requests came before the banner was visible.} An ad pixel `fired' when a Shopify sandbox sent an ad request before any choice (Default load), as in Figs~3 and 4. Colours, one meaning each: brick, consent lost (the first sandboxed requests in (b); the region default that allows them in (d)); slate, the comparison (d); amber, the banner becoming visible; dark slate, Shopify's sandbox (c); black and grey, stores with and without a banner (a, e), a contrast that is not about consent; other greys, context. (a)~Stores on which an ad pixel of each vendor group fired, among the 409 stores that showed a banner (black circles) and the 949 that did not (grey squares), each named on the `Any vendor' row, which is set apart below. Horizontal lines in (a) and (d): Wilson 95\% confidence intervals (CIs), in the colour of their mark. Keys in (c) to (e) sit above their plots. (b)~The 205 banner stores with a pre-choice sandboxed request: cumulative share whose first sandboxed ad request had been sent by time $t$, timed from the banner becoming visible (amber line; the time axis of Fig.~1b,c); thin light grey line: one polling interval of the banner detector (250~ms) before the banner. The pale brick ground under the curve, left of the amber line, marks the stores that sent before the banner: the open ring where the curve meets the amber line is 151 of 205 stores (73.7\%), and the dot on the light grey line the 108 (52.7\%) that sent more than one poll before it. The 9 stores whose first request lies beyond the axis are counted at its ends, so the curve starts above the bottom of the axis and ends below its top; where it leaves the axis it ends in an arrowhead. In the 46 such stores loaded twice, `request first' agreed in 36 ($\kappa$ 0.54). (c)~Origin of all ad requests in the Default load, by vendor group: Shopify sandbox (dark slate, the colour of the sandbox frame in Fig.~1a; web worker or sandboxed iframe, including frames nested in it), Shopify's open pixel (pixel type `OPEN', which runs in the page; mid grey) and ordinary page scripts and frames (light grey); thin white gaps separate the segments. Requests relayed from a sandbox by Shopify code in the page are Google's only (115 of 20,503; 0.6\%): they keep their share of the bar but are narrower than a segment gap and are not drawn. Sandboxed share of each group's ad requests: Google 20.1\% (4,112 of 20,503), TikTok 55.1\% (2,280 of 4,135), Meta 16.5\% (408 of 2,478), Pinterest 62.2\% (1,189 of 1,911), Snap 31.0\% (612 of 1,973), Microsoft 46.3\% (1,093 of 2,363) and Reddit/other 43.7\% (2,910 of 6,653); across all vendors 31.5\% (store-bootstrap 95\% CI 29.8--33.2) of 40,016 requests were sandboxed. (d)~Banner stores on which any ad pixel fired, within five consent products, by Shopify's region default: `allowed until choice' (brick circles) and `denied until consent' (slate squares). These are the shares of Fig.~3a's four product rows, plus Shopify's banner, which Fig.~3a lists in its caption but does not draw, here under Fig.~3a's axis title with the stores under each default: Pandectes 49 and 44; OneTrust 68 and 21; Cookiebot 37 and 8; Osano 7 and 9; Shopify's banner 1 and 65; grey marks: fewer than 10 stores under one default (Cookiebot `denied until consent', both of Osano's, whose name is grey, and Shopify's banner `allowed until choice'). The pooled Mantel--Haenszel difference over 12 products is in Fig.~3a (right). (e)~How many sandboxed ad requests each store sent before any choice (Default load, no click): share of stores with at least $n$ requests (complementary cumulative distribution; logarithmic $n$ axis), for the stores that showed a consent banner (black) and those that showed none (grey). Marks at the left end, at least one request: 205 of 409 banner stores and 598 of 949 stores without a banner; Wilson 95\% CIs 45.3--54.9\% and 59.9--66.0\%; a cluster bootstrap by consent tool widens them to 32.6--58.9\% and 56.8--68.0\%.}
\label{fig:s6}

% ---- Figure S7 (figS7.pdf)
\caption{\textbf{Blocked on paper in the main sample: vendors that a blocklist named still fired from Shopify's sandbox in every subset, while the blocker acted only on the page.} Colours: brick, sent before any choice (consent lost); slate, the comparison route (a page script); greys, context. All counts are for the main sample (Tranco top 200,000 ranks); the expansion to ranks 200k--1M and the combined estimate (primary set: 46 of 89 stores) are in Fig.~4a. (a)~Blocklist stores on which the vendor that the consent-management platform's (CMP's) blocklist declared blocked fired: its own sandboxed Shopify pixel sent an ad request before any choice (Default load). Top, set apart, the primary estimate (filled brick circle): stores whose CMP wrote no `marketing: yes', 50.0\% (7 of 14). Below, the sensitivity subsets (open brick circles): all 21 blocklist stores (12 fired); the stores with the vendor's sandboxed pixel configured (10 of 15); without the 4 stores whose CMP wrote `yes' by itself (2 of them because Pandectes' bot rule detected our automated browser; 9 of 17); and without the 5 stores with a `yes' before the pixel started (9 of 16). Horizontal lines: Wilson 95\% confidence intervals (CIs). (b)~The blocklist entries on the 21 blocklist stores. Left: the 29 entries whose vendor has a configured sandboxed Shopify pixel (15 stores) as a 2~$\times$~2 matrix of counts: whether the vendor sent an ad request before any choice from Shopify's sandbox (rows) and, counted independently, from an ordinary page script (columns); both routes 6; the sandbox only 15 (pale brick cell: consent lost only through the sandbox; the cells are split by two crossing rules, as in Fig.~S4a--c); a page script only 2; neither 6. Right: share of entries sent before any choice, with Wilson 95\% CIs: the same entries from the sandbox (brick circle; 21 of 29) and from a page script (slate square; 8 of 29) and, set apart, the 75 entries without an identified configured sandboxed pixel (20 stores; classic tag, open pixel or unknown installation) by any route (grey circle; 15 of 75). (c)~One re-check load per blocklist store, after the main crawl, drawn as the ledger of Fig.~2a (each box: the stage, left, and its number of stores, right; the stores that leave a stage branch right, with their reason in grey, so that each count is the one above minus its branch; the brick arrow marks the leak). Of the 21 blocklist stores the blocker was active in 18; on 8 it neutralised the tag of a vendor it declared blocked; on 5 of these that vendor's sandboxed pixel still sent an ad request. (d)~The primary set store by store: its 19 blocklist entries whose vendor has a configured sandboxed Shopify pixel, on 11 stores, as two store $\times$ vendor-group matrices with one layout, each on the ground of its region in Fig.~1a. Left, on the darker grey ground of the pixel sandboxes: the route through Shopify's sandbox (12 of 19 entries sent an ad request before any choice); right, on the light grey ground of the page: the route through an ordinary page script (5 of 19). Columns are the stores, grouped by the consent tool whose blocklist declared the vendor blocked (named below a rule under its columns); rows are the vendor groups used in every other figure (the one Reddit/other entry is Amazon Ads); brick: sent by that route; grey: not sent by that route (key above the right matrix); blank: no entry. The same position in both matrices is the same entry: 4 entries were sent by both routes and 6 by neither. Blocklists are readable for Pandectes, OneTrust and Cookiebot only.}
\label{fig:s7}

% ---- Figure S8 (figS8.pdf)
\caption{\textbf{After a choice, new sandboxes kept sending ad requests, and where a consent tool wrote `no' without a click a sandbox had almost always started first.} Colours: brick, consent lost; slate, the comparison; amber, the consent tool's write; greys, context. Every key sits above its plot. (a)~The per-vendor estimates of Fig.~5a, repeated here (they are not new estimates) as rows on the scale of (b), so that the vendors can be read against the subgroups of (b): on the 231 stores where the Decline click worked, stores on which a new sandbox sent an ad request to each vendor group under Default (filled black circles), on the fresh page after Decline (open squares; new sandboxes, 10~s) and under Global Privacy Control (GPC; filled grey triangles; first 10~s); `Any vendor' is set apart below. Lines in (a) and (b): Wilson 95\% confidence intervals (CIs), in the colour of their mark. (b)~The fresh-page measure of (a), any vendor (open squares), by banner type (top three rows; each grouping is named in grey above its rows): Shopify's own banner 9 of 57; a consent-management platform (CMP) that writes to Shopify 32 of 143; a blocker only, unknown or none 12 of 31; and by the state Shopify recorded after the click (bottom three rows): explicit refusal 37 of 196; `marketing' still allowed 13 of 20; no record 3 of 15. (c)~What the pixels that still sent requests declare, by window after a Decline (same page after the click; unload beacon at the reload; fresh page) among the 196 stores where Shopify recorded a refusal, and under GPC among the 1,354 stores where Shopify set `sale of data: no'. Share of the window's stores whose still-sending pixel declares a purpose the visitor refused (brick), only purposes the choice left allowed (grey) or no purpose (dark grey), in the class names, colours, key order and bar form of Fig.~5b: one thin bar per class from zero, side by side rather than stacked because a store can be in several classes, with its Wilson 95\% CI as a dark line, as in Fig.~5b; a class with no store has no bar, only its interval. Stores whose pixel declares a refused purpose: 19 (same page), 20 (unload beacon), 0 (fresh page) and 0 (GPC); only purposes left allowed: 3 (same page), 3 (unload beacon), 13 (fresh page) and 111 (GPC); no purpose: 8 (same page), 6 (unload beacon), 27 (fresh page) and 191 (GPC). A store can appear in several classes and windows (one class per store: Fig.~5b). (d)~The race in Fig.~5c's raster and time axis: one row per page load in which a CMP wrote `marketing: no' or `yes' without a click, in four blocks (`no' in the Default load, 59 loads; `no' under GPC, 87; `yes' in the Default load, 22; `yes' under GPC, 13), each block sorted by its loads' first third-party sandbox that declares `marketing'. One tick per load, at that sandbox's start minus the CMP's consent write (amber line, named above the plot), on Fig.~5c's symmetric-log axis (linear close to the write, logarithmic beyond; the two innermost tick labels of Fig.~5c are omitted here). Brick: the sandbox started before the `no'; grey: after the `no'; slate: after a `yes'. A wider gap sets the two `yes' blocks apart from the two `no' blocks. The pale brick ground between the brick ticks and the amber line marks the time each of those sandboxes ran before the `no'. In 144 of 146 `no' loads a sandbox started first (the other 2 are GPC loads); in all 35 `yes' loads it started after. Stores where a CMP wrote `marketing: no' without a click and a sandboxed ad request was still sent before any choice: 39 of 65 (60.0\%; Wilson 95\% CI 47.9--71.0\%). The per-sandbox version is Fig.~5c.}
\label{fig:s8}

% ---- Figure S9 (figS9.pdf)
\caption{\textbf{Vendors received the full landing URL in every load, refusal signals mostly in Google's requests, and a creator's per-video code on every load of Store C.} Colours: brick, requests that carried the per-video code; slate, Google beacons whose Consent Mode state denied ad storage; amber, the banner becoming visible and the Decline click; greys, context. (a)~Sandboxed ad beacons carrying the full landing URL with our test code, by vendor group and load: Default (filled black circles), Decline (open squares) and Global Privacy Control (GPC; filled grey triangles); key above the plot; Korn--Graubard 95\% confidence intervals (CIs) clustered by store. Stores per load (Default / Decline / GPC): Google 390 / 393 / 172; TikTok 346 / 344 / 16; Meta 18 / 18 / 4; Pinterest 264 / 264 / 14; Snap 179 / 178 / 6; Microsoft 226 / 224 / 64; Reddit/other 314 / 309 / 84; all vendors 769 / 764 / 242. A load with fewer than 10 stores for a vendor (Meta and Snap under GPC) is not drawn. Markers in (b) and (c) as in (a). (b)~One plot on the scale of (a), in two groupings named in grey above their rows. `Ad requests': ad requests carrying a refusal signal in the URL (Google: Consent Mode, `npa', `rdp' or `us\_privacy'; Microsoft: `asc=D'), in the Default load, on the next page after Decline and in the GPC load; 95\% CIs from resampling stores (percentile bootstrap, 2,000 resamples). On the fresh page after Decline (open squares): Google 434 of 1,220 (35.6\%), Microsoft 80 of 408 (19.6\%), all other vendors 0 of 1,211. `Stores, all flagged': stores still sending sandboxed ad requests on which every one of them carried a refusal signal (open square: the fresh page after Decline; grey triangle: GPC): 9 of 53 on the fresh page after Decline and 18 of 279 under GPC (Wilson 95\% CIs). (c)~Google's ad requests only. Left: beacons by the Consent Mode state in their `gcs' parameter: ad storage denied (G100 or G101; slate), granted (G111; dark grey) or no state (light grey); thin white gaps separate the segments; key above the plot. Beacons: 20,480 in the Default load, 1,219 on the fresh page after Decline and 16,703 in the GPC load; other states (at most 0.3\% of a load's beacons) are too few to draw. Right: share of Google's ad requests carrying each signal (`gcs' or `gcd' with ad storage denied, `npa=1', `rdp=1', a `us\_privacy' opt-out; markers as in (a)). These are request-level shares of the same requests (Default 20,503; Decline 1,220; GPC 16,718); because requests cluster within stores, their 95\% CIs come from resampling stores, as in (b); an interval narrower than its marker is hidden behind it. (d)~Store~C, a creator-linked store, reached from the creator's video link: every ad request in five loads (Default and its repeat, Decline, GPC and its repeat) as an event strip on one broken time axis with one scale. The axis is cut (double slashes) where nothing was sent or happened for a long span, and two short, sparse stretches share the middle stretch. Each load has a lane of requests from a sandbox (`Sandbox', on the darker grey ground of Fig.~1a's pixel sandboxes) and a lane of all others (`Other', on the light grey ground of the page). One tick per group of requests of one vendor group and kind less than 50~ms apart, at their mid time (37 ticks stand for more than one request); tall ticks: data beacons and other requests, brick when at least one request of the tick carried the per-video code and grey when none did; short light grey ticks: library loads, none of which carried the code (key above the plot). Amber lines: the banner becoming visible and the Decline click; grey line: the reload in the Decline load; each event is named above its own line, the banner once. In the sandbox lanes of the Default, repeat and Decline loads the first two ticks are Google's, the two short light grey ticks are TikTok's library loads and the last tick stands for the 4 code-carrying TikTok requests counted in the caption of Fig.~6a; nothing was sent from a sandbox in the GPC loads. The code reached 5 vendors (3 other than Google) in the Default load and 5 (3) in its repeat; in the Decline load 5 before the click and 4 after the reload (0 from a sandbox); under GPC 3 (1) and 3 (1) in its repeat. Signals carried in request bodies are not visible.}
\label{fig:s9}

% ---- Figure S10 (figS10.pdf)
\caption{\textbf{Beyond Shopify: in a pilot of selected Wix, BigCommerce and first-party-gateway sites, ad requests were common before any choice, after a working Decline and under Global Privacy Control (GPC).} The sites are a pilot of selected sites, not a sample: 80 Wix and 80 BigCommerce sites (80 and 77 captured) and 26 sites whose Google tag loads through a first-party tag gateway or tagging server (`first-party gateway'; 24 captured). They were crawled on 27 Sep 2026 from one California vantage point with the same three loads as the main study (Default, Decline, GPC). An ad pixel `fired' when the site sent an ad request: a request to an advertising vendor, or a first-party request whose wire format names Google Ads, Meta or TikTok (the axis titles of Figs~3 and 4). Forwarding onward from a server (server-side) is not observable and is never inferred. The three setups are the rows of every plot. (a)~Banner sites on which any ad pixel fired in the Default load, before any choice: Wix 28 of 57; BigCommerce 48 of 72; first-party gateway 17 of 19. (b)~Sites where the Decline click worked. Left: an ad pixel fired after the click, on the same page or the next page (Wix 14 of 33; BigCommerce 32 of 53; gateway 7 of 8). Right: a first-party collect beacon after the click, that is a GA4, Google Ads, Universal Analytics, Meta or TikTok beacon sent to the site's own domain or tagging server: Wix 5 of 33; BigCommerce 45 of 53; gateway 6 of 8. (c)~All sites with a Default and a GPC load: any ad pixel fired in each load, Default above GPC (Wix 40 and 40 of 80; BigCommerce 52 and 48 of 77; gateway 22 and 20 of 24). Markers, in one key above (c): Default load (filled black circle), Decline load (open square), GPC load (filled grey triangle); lines: Wilson 95\% confidence intervals; light grey marks: fewer than 10 sites (the gateway sites with a working Decline click, in (b)); every plot has the same scale. Two case studies give no rates and are not drawn. Partytown, which runs third-party scripts in a web worker: of 8 selected sites, 5 were captured; requests from Partytown workers numbered 0 in all loads, and the 3 sites with an ad request in the Default load sent it from the top frame's main thread (3 sites) and from iframes (1 site). Zaraz, which runs tools on Cloudflare's servers: on the 1 selected site Zaraz loaded, made 0 track calls in all loads, and sent no ad request in any load. Rates describe the selected sites only: they are not population estimates, the setups are not compared statistically, and they are not comparable with the Shopify rates of the main study, which count ad requests from Shopify's sandbox.}
\label{fig:s10}
